\section{What ``bounded over the primes'' precisely means}\label{sec:primes}
The following arithmetic form supplies the missing qualification in the sketch's prime-language formulation. Closely related bounded-complexity lifting theorems are \citet[Definition~2.3 and Theorems~3.1, 3.4]{BrunjesSerpe09}. The form proved here allows the clearing multipliers to vary without an a priori bound, provided the degrees are those of the \emph{integral} effective parts and the realization is a fixed global section.

\begin{theorem}[P6: arithmetic lifting with varying multipliers]\label{P6}
Let $F\hookrightarrow\C$ be a number field, $U\subset\operatorname{Spec}\cO_F$ a nonempty open with $\ell$ invertible, and $f:\mathcal X\to U$ smooth and projective. Fix a relative ample $L$ and a nonzero section
\[
 t_\ell\in H^0(U,R^{2p}f_*\Q_\ell(p))
\]
whose geometric generic value corresponds to a rational Betti class $t_B$ under the chosen comparison. The following conditions are equivalent:
\begin{enumerate}[label=(\roman*)]
\item $t_B$ is rational algebraic;
\item some fixed positive $N$ and finite $D$ give integral effective representatives of $Nt_\ell$ with both degrees at most $D$ at every geometric residue fibre outside finitely many places;
\item the same holds at infinitely many places, with $N,D$ fixed;
\item at infinitely many places it holds with $D$ fixed and positive integers $N_v$ allowed to vary without any assumed bound.
\end{enumerate}
In (iv), one of the occurring integers $N_0$ already gives a characteristic-zero effective pair with the same degree bound $D$. No bound on the fields of definition of the geometric residue cycles is required.
\end{theorem}

\begin{proof}
We first specify the parameter-space fact needed in mixed characteristic. After replacing $L$ by a fixed relatively very ample power, the bounded-degree relative cycle presheaves are $h$-represented by projective Chow schemes over the noetherian base; this is \citet[Theorem~4.4.11 and Corollary~4.4.13, manuscript pp.~79--81]{SuslinVoevodsky00}. Take the finite union of their products for the allowed positive and negative degree pairs. By the definition of sheafification, a finite type $h$-cover represents the universal sections by genuine effective integral relative cycles. Quasicompactness permits a finite covering family and a common refinement for both cycles.

On each of the finitely many reduced irreducible parameter components, simultaneously flatten the closures of the finitely many generic supports. Proper modifications are surjective, and the strict-transform families are flat and projective. Their associated relative cycles equal the pulled-back original cycles because relative cycles are determined by their generic values. This uses the flattening theorem in the same source, Theorem~2.2.2, and its Section~3.2 construction of cycles of flat subschemes; the detailed proper-cover form is also given by \citet[Lemmas~6.3.19, 6.3.24 and Theorem~6.3.25, pp.~188--193]{Anderson19}. We have obtained a finite type $U$-scheme $T$, surjective on geometric cycle pairs, with actual flat cycle families. Degrees remain the prescribed degrees. We do not assert the existence of a universal flat \emph{subscheme} on the coarse Chow scheme.

The finitely many irreducible components of $T$ which do not dominate $U$ lie over a finite set of closed places $E$. Choose one of the places in (iv) outside $E$, and a parameter point $y$ representing its pair. Choose an irreducible component of $T$ through $y$; it is horizontal. A discrete valuation ring maps its generic point to the generic point of this component and its closed point to $y$, allowing residue-field extension. This is the trait realization of specialization, Stacks Project, Lemma~28.5.10, Tag~054F \citep{Stacks054F}. Pull back the flat cycle families to this trait. Their generic and special cycles are effective and have the same degrees.

Smooth proper base change gives an isomorphism
\[
 \mathrm{sp}:H^{2p}(X_{\bar K},\Q_\ell(p))
 \xrightarrow{\sim}H^{2p}(X_{\bar\kappa},\Q_\ell(p)).
\]
The exact Grothendieck locators are SGA~4, Expos\'e XVI, Corollaire~2.2, for this specialization, and SGA~6, Expos\'e X, Appendix~\S7, equation~(7.13.10), original pp.~578--579, for its compatibility with rational cycle classes \citep{SGA4,SGA6}. Hence
\[
 \mathrm{sp}\bigl(\cl(Z_K^+-Z_K^-)-N_vt_{\ell,\bar K}\bigr)=0.
\]
Injectivity gives the generic equality with that \emph{same} positive integer $N_v$. The fraction field is finitely generated over $F$ and embeds in $\C$ extending the chosen embedding of $F$. Betti--\'{e}tale comparison for cycle classes gives
\[
 \cl_B(Z_\C^+-Z_\C^-)=N_vt_B.
\]
Both sides are rational Betti classes, so equality after scalar extension to $\Q_\ell$ is equality over $\Q$. This proves (iv)$\Rightarrow$(i). Specializing the characteristic-zero parameter component to a $\bar F$-point preserves its cycle class and degrees. Thus one can take the pair over a finite extension of $F$, with $N_0=N_v$ and the original $D$.

Conversely, a rational cycle for $t_B$ can be spread from $\C$ to a finite type $\bar F$-parameter scheme and then specialized to a closed point of a connected component. Cycle-class constancy gives a cycle over $\bar F$. Clear its coefficients and separate its two signs. Over a finite extension of $F$ all components are defined; after deleting finitely many places their closures are flat. Their geometric residue cycles have the original degrees and specialize to $Nt_\ell$. This proves (i)$\Rightarrow$(ii). The remaining implications are immediate.
\end{proof}

One can also see explicitly why varying $N_v$ is harmless. On a fixed horizontal parameter component, an equality with $N_v$ transports through a trait to its geometric generic cycle class. Two such integers would give $(N_v-N_w)t_\ell=0$ there; since the section is nonzero, they are equal. Only finitely many components occur. Thus outside finitely many places the possible multipliers for degree at most $D$ belong to a finite set. This observation does not compute their maximum.

\paragraph{Exact application of 032.}
The supplied \citet[Remark~8.5, printed pp.~50--51]{OpenAI26} concerns an abelian variety over $\overline\Q$ that need \emph{not} be CM. It states that every rational Hodge class specializes, at every good prime, to one rational algebraic cycle compatible simultaneously with the prime-to-$p$ and crystalline realizations. The remark uses the CM transport cited there as [32], Corollaries~3.4 and~6.4. With the compatible global realization fixed over a sufficiently large field of definition, it supplies pointwise cycles of the kind in P6. It supplies no common degree bound after their rational coefficients are cleared. P6 therefore makes the arithmetic gate exact: uniform integral effective degree at infinitely many good places is equivalent to characteristic-zero algebraicity. Pointwise algebraicity in all good characteristics alone is not the lifting argument.
